\chapter{The Beginning}

The start of computational theory can be traced back to the 1930s, when mathematicians and logicians were trying to formalize the notion of computation and algorithms. The Church-Turing thesis, proposed independently by Alonzo Church and Alan Turing, laid the foundation for the study of computational complexity.

\begin{theorem}[Church-Turing Thesis]
	Every effectively calculable function is computable by a Turing machine.
\end{theorem}

Then mathematicians and computer scientists began to explore the limits of computation, with such an open thesis:

\begin{problem}[Extended (Quantum) Church-Turing Thesis]
	Every physically realizable model of computation can be simulated by a (quantum) Turing machine efficiently.
\end{problem}

\section{Definition of Computational Theory}

\begin{remark}
	Computational theory deals with how much resources we need to solve a computational problem.
\end{remark}

Resources can be many things. For example:
\begin{itemize}
	\item Time
	\item Memory / Space
	\item Samples Complexity (how many samples we need to distinguish between distributions)
	\item Communication Rounds
	\item Oracle Calls
	\item Randomness
	\item Power
	\item Quantum Entanglement
	\item Number of Parallel Workers
	\item $\cdots$
\end{itemize}

\section{What is Complexity About?}

In CS Theory, people do research in algorithms and complexity. People might lean on one side or the other, but they are closely related. In a word:
\begin{remark}
	Algorithms are about the upper bound of a problem, while complexity is about the lower bound of a problem.
\end{remark}

\chapter{Time and Non-deterministic Time}

\section{Languages and Problems}
In computational complexity, we talk about languages as how we think about computational problems.

\begin{definition}
	A \textbf{language} $L$ is a set of strings that $L\subseteq \Sigma^*$, where $\Sigma$ is a finite alphabet (often $\{0,1\}$ bits.)
\end{definition}

\begin{eg}
	An algorithm $\text{PRIMES}(n)$ can decide whether a number $n$ is prime or not. The language of this problem is:
	\[ \text{PRIMES} = \{ \langle x \rangle : x \text{ is prime} \}. \]
	$\langle x \rangle$ is the encoding of $x$ in binary.
\end{eg}

\section{Classes of Complexity}

\subsection{Deterministic Time Complexity}
\begin{definition}
	$L \in : \DTIME(T(n))$ if there's a deterministic Turing machine $M$ that decides $L$ in time $O(T(n))$, where $n$ is the size of the input.
	$ \P := \bigcup_{c\ge 1} \DTIME(n^c)$ is the class of problems that can be solved in polynomial time.
\end{definition}

Researchers have found a polynomial time algorithm for the PRIMES problem in 2002, so $\text{PRIMES}\in \P$.
